\selectlanguage{english}
\Language{en}{English}{User guide}
\firsttopic{Open, sleep, return}
Inkline runs a local shell on reMarkable 2. Use a Type Folio or USB keyboard,
work on the tablet, or connect with SSH or Goblin Mosh. Text and graphics are
rendered in grayscale.

\begin{notice}
\key{Opt+RightAlt+T} opens Inkline from the tablet itself.\par
\key{Power button:} press and release to suspend; press again to wake.
Your open shells and editor buffers stay in RAM.
\end{notice}
Sleep pauses the tablet; it does not exit Inkline or save settings to flash.
Network connections may need to reconnect after waking. The wake press is
ignored as a sleep request, so the tablet does not immediately suspend again.
After an update, quit and reopen Inkline once to activate the new power handling.

To return to notebooks, tap \key{Quit} or press \key{Ctrl+Shift+Q}, then confirm.
Typing \code{exit} closes one terminal; closing the last returns to notebooks.
The bottom bar has Escape, battery status, Quit and Settings.
The startup guide and tiny Goblin logo disappear with \code{clear}.

If a shortcut app owns the screen, \key{Opt+RightAlt+T} returns to Inkline.
For emergency recovery, hold \key{Opt+RightAlt+Backspace for two seconds}.
Recovery closes every Inkline terminal: unsaved work in those sessions is lost.
The notebook interface remains the default after reboot.

RightAlt means the Alt/Opt key to the right of the spacebar. Hold it with the separate Opt key, then tap T; do not hold Ctrl. Plain Super remains available for personal bindings.

\ProjectLinks

\topic{Keys and terminal slots}
The Type Folio has a separate \key{Opt} key between Ctrl and Alt, and a
\key{right Alt/Opt} key. They have different jobs.

\begin{keytable}
\key{Keys} & \key{Action} \\ \midrule
Separate Opt + 1--0 & F1--F10 \\
Right Alt/Opt + 1--9 & Select terminal slot 1--9 \\
Right Alt/Opt + Left / Right & Previous / next slot \\
Right Alt/Opt + Tab & Escape \\
Right Alt/Opt + Up / Down & PageUp / PageDown \\
Right Alt/Opt + Backspace & Confirm quitting all terminals \\
Right Alt/Opt + Space & Settings \\
Left Alt + Space & Unicode keyboard \\
Right Alt/Opt + C / V & Copy selection / paste \\
\end{keytable}

Ctrl+Shift+B passes through to terminal programs, including Emacs.

On the \key{US Type Folio}, right Alt/Opt+\key{0} types \code{+}.
Right Alt/Opt with the \key{minus key just left of Backspace} types \code{=}.
With the input method off, Shift+6 types a literal caret: \code{c^2} stays
\code{c^2}. Keyboard zoom shortcuts are disabled.

USB keyboards keep their ordinary function keys. The separate Opt shortcut
uses Meta when a USB keyboard provides it.

There are at most \key{nine independent slots}. Selecting an empty slot opens
a shell. The indicator counts open terminals: with only slots 1 and 9 open,
slot 9 shows \key{Terminal 2 / 2 · slot 9}. Closing a shell frees its slot
without renumbering the shortcuts.

Each terminal has its own shell, directory, input composition, graphics and
500 lines of scrollback. Display and input-method settings apply to all slots.

\topic{Settings and display}
Tap the rightmost \key{Settings} button or press \key{right Alt/Opt+Space},
including while the bar is hidden. Use \key{left Alt+Space} for Unicode.
A filled choice is active; a dashed outline is keyboard focus. Tab moves
focus; arrows and Enter change choices.

\key{Caps Lock:} choose Control (the default) or Caps Lock.
\key{Text size:} pinch with two fingers, or use the minus/plus buttons in
Settings. Slow movement changes the size one pixel at a time, from
\key{6 to 48 pixels}.

\key{Text darkness:} 50\% is the original rendering. Higher values darken
letter edges. \key{Minimum contrast} separates foreground and background
luminance when an app chooses similar colors; the default is 35\%.

\begin{choices}
\key{E-paper mode} & \key{Tradeoff} \\ \midrule
Fast & Minimal batching; animation waveform; more ghosting \\
Balanced & UI waveform; 32 ms output batching \\
Crisp (default) & Cleaner grayscale; content waveform; slower \\
Mono & Fast black and white; discards image grays \\
Saver & 120 ms batching; fewer updates during bursts \\
\end{choices}

Only changed text rows are redrawn. The physical display still takes time
to update. Saver reduces update requests; battery savings have not been measured.
An update preserves an explicitly saved mode instead of forcing the new default.

\begin{notice}
All settings stay in \key{RAM} during use and are saved together when Inkline
exits normally. Finger lifts, closing Settings and sleep do not write them
to flash. An unchanged session writes nothing. A crash, forced kill or power
loss can lose unsaved setting changes.
\end{notice}

The bottom bar reads battery state once per minute from Linux's read-only
power-supply interface. Run \code{inkline-battery} for percentage and charging
state, \code{inkline-battery --percentage} for one compact value, or
\code{~/inkline battery}. The installer uses the compact value in Inkline's
no-color Bash prompt without changing prompts outside Inkline.

\topic{Unicode and input methods}
\key{Left Alt+Space} opens the Unicode keyboard. Tap a character to insert it;
the keyboard stays open. Tap \key{Done} or press Escape to close it.

Categories include letters, marks, numbers, punctuation, math/currency,
symbols, spaces, format characters and private use. Drag vertically to
scroll; arrows, PageUp/PageDown, Home/End and the mouse wheel also work.
\key{Tab} changes category; Enter inserts the selected character.

Type a hexadecimal Unicode value, such as \code{03bb}, then Enter to insert
λ. Backspace edits the value. Combining marks show a dotted circle as a
visual aid; only the mark is inserted. Spaces and format characters have labels.
The catalog uses Qt's Unicode tables, excluding controls, unassigned values,
surrogates and noncharacters. Inkline bundles OFL-licensed Noto CJK and Symbols
fonts, plus GNU Unifont and Unifont Upper as broad last-resort fallbacks. No
practical single font covers every Unicode glyph, and private-use values still
require the font that defines them.

Press \key{F2} in the picker or tap \key{Settings}, then choose \key{Input method}:

\begin{choices}
Off — direct keyboard & Ordinary US typing \\
Romaji & Japanese hiragana from romanized input \\
Pinyin & Chinese from pinyin \\
Zhuyin & Chinese with the basic Zhuyin map \\
Wubi 86 & Chinese shape codes \\
US-International & Common Latin accents \\
\end{choices}

Choose a result in the candidate strip. Select \key{Off — direct keyboard}
in the same menu to return to normal US input. The choice applies to all
terminals and is saved when Inkline exits.

\topic{Touch, pen and clipboard}
\key{Two fingers up/down:} scroll history. Drag down for older output and
up toward the prompt. Each terminal retains 500 physical lines in RAM,
in addition to the live screen.

\key{Two fingers left/right:} switch one terminal slot per gesture.
Lift your fingers before switching again. Spread or close them to resize text.

\key{One finger:} drag to select text. Right Alt/Opt+C copies;
right Alt/Opt+V pastes. All terminals share one RAM clipboard, up to 128 KiB.

\key{Pen:} apps that enable terminal mouse reporting receive pen press,
movement and release as left-mouse events, using the negotiated protocol
(including SGR). Otherwise, dragging selects text. Hold \key{Shift} to force
local selection in an app that uses mouse reporting.

\key{OSC 52} lets terminal programs, including remote programs over SSH,
read or replace Inkline's RAM clipboard. It is separate from the notebook
clipboard. Right Alt/Opt+X copies and explains that deletion belongs in the
editor: OSC 52 cannot delete an editor's text.

\key{Graphics:} kitty inline and shared-memory transport stay in RAM;
file and temporary-file transport are disabled. Sixel works when requested
explicitly, but is not advertised automatically. Kitty placeholders,
animation playback and some sixel modes remain unfinished.

\code{clear} and full-screen erase remove visible graphics as well as text.
Off-screen graphics are reclaimed under memory pressure. There is no disk
image cache; large images still consume the tablet's limited RAM.

\topic{USB keyboard and typewriter}
Open \key{Settings → USB (page 3)}. \key{Network} restores USB Ethernet;
\key{Send keyboard} makes the tablet a keyboard for another computer;
\key{Receive keyboard} accepts a wired keyboard through a powered OTG adapter.
Type Folio stays available in every mode. Install the website's Python utility
for outgoing keyboard mode; no extra Python modules are needed.

Choose Send keyboard, connect USB-C to the receiving computer, select a host
profile below, then tap \key{Open typewriter}. Live output starts automatically
when Send keyboard is active. When output is paused, the editor displays
a black \key{PAUSED · USB OUTPUT OFF} badge; live output has a separate outlined badge.
\key{Escape} or \key{Stop} pauses and releases keys. Leaving the typewriter,
switching terminals, losing window focus, opening a native app or sleeping
also pauses output. Nothing is replayed automatically after reconnection.

\begin{choices}
US / ASCII & US keyboard layout, Caps Lock off. Non-ASCII text is rejected. \\
Mac & Select Unicode Hex Input in Keyboard → Input Sources; Caps Lock off.
Characters through U+FFFF only; supplementary characters are rejected. \\
Linux & US layout, Caps Lock off. The destination app/input method must accept
Ctrl+Shift+U, hexadecimal digits, Enter. This does not work in every app. \\
Windows & US layout, Caps Lock off. Install WinCompose on the PC, use Right Alt
as Compose and enable Unicode input. Inkline sends Compose, u, six hex digits, Enter. \\
\end{choices}

Get WinCompose from \url{https://wincompose.info/}. Custom compose rules must
not override the Unicode sequence. There is no universal HID Unicode text
channel: host setup and fonts determine the result. Test a short sample first.

\key{Input method} uses the same engine as the terminal: Off/direct US,
Romaji, Pinyin, Zhuyin, Wubi 86 or US-International. Only committed characters
leave the tablet. Changing the method pauses typing; tap Start again.
Choose \key{Off} to return to normal US input. \key{Unicode} or left Alt+Space
opens the character picker and sends selected characters to the USB host.

Typewriter is a local editor whose document stays in RAM while Inkline runs.
Ordinary typing changes that document and, while live typing is on, sends each
committed character to the receiving computer. Copy, cut, paste, cursor movement,
touch selection and two-finger scrolling act on the local editor.

\key{Files \& send} accepts any UTF-8 filename; Load and Save do not require a
particular extension. Select 5, 10, 20, 40 or 80 characters per second, then
send the entire editor. \key{Stop sending document} cancels a bulk send.
Editing does not write the document to storage. Save writes it explicitly;
normal Inkline exit writes the current named file or uses the hidden
\path{~/.inkline-typewriter-draft} recovery name for an untitled document.
\key{Exit typewriter} returns to USB Settings.

\topic{Send a file; restore networking}
The installer adds \code{inkline-usb} and \code{keyboard-send} to
\path{~/.local/bin}. \code{~/inkline usb}, \code{~/inkline type}, and
\code{inkline-type} remain compatibility aliases. The separate Outpost tools
remain installed.

\UsbExamples

UTF-8 files have \key{no fixed size limit}. The sender reads small chunks,
without storing the document in RAM or creating a spool file. Ordinary,
seekable files are checked completely before any keys are sent, then read
again; keep the source unchanged while sending. A pipe is validated as it
arrives, so an error can follow already-sent text. An optional UTF-8 BOM is
removed and CRLF/CR become newlines. Enter and Tab act on the focused app.

\code{--cps} (default 20) and \code{--delay-ms} select typing speed;
Unicode needs more reports per character. \code{--start-delay} gives time
to focus the destination. Dry-run prints reports without using USB.
Ctrl+C and \code{--stop} cancel the sender and release keys. Only one sender
can run at once. The interactive queue is bounded; overflow pauses typing.
After interruption, inspect the destination before retrying a partial send.
Text, runtime files and helper caches do not write to flash; the host profile
is saved with other settings on normal exit.

Tap \key{Network} or run \code{inkline-usb network} to restore USB Ethernet.
Leaving Inkline restores networking, including after a crash or emergency
restart. Mode changes are not saved for boot. To schedule independent recovery:
\begin{shell}
inkline-usb recover-after 90
inkline-usb send-keyboard
\end{shell}
Switch away from networking on the tablet or through Wi-Fi SSH, not USB SSH.
A failed transition attempts Ethernet recovery. Inkline refuses to take USB
from an active satellite PPP session or an unknown gadget; this also applies
to exit cleanup. The emergency Opt+RightAlt+Backspace chord still works,
but restarting Inkline closes its terminals and loses unsaved work.

\topic{Launch programs with keys}
Register shortcuts in an Inkline shell. Arguments are literal; use an explicit
shell for shell syntax. Shortcut programs read \code{~/.bashrc} through Bash.

Settings → Command shortcuts (page 2) lets you assign a command to a letter. Choose Terminal or Native e-paper app, then Save. Remove asks for confirmation; Reset draft discards edits. T is locked. Only Save and confirmed Remove write to storage. The command-line tool manages the same bindings.

All global shortcuts now use Opt+RightAlt, including T and the held Backspace recovery chord. Use the Folio’s separate Opt key; on USB keyboards use Windows/Command (Super). Ordinary Ctrl+Alt remains available to Emacs. Quit and reopen Inkline after upgrading; installation preserves open terminals.

\begin{shell}
~/inkline shortcut register e emacs
~/inkline shortcut register p python3
~/inkline shortcut list
~/inkline shortcut deregister p
\end{shell}

\key{Opt+RightAlt+E} opens Emacs in a free slot.
\code{~/inkline run PROGRAM ARG...} launches without a binding.
Registering replaces a binding; deregistering leaves running programs alone.
Keys use US physical positions. Quote punctuation and use full paths for
custom programs; programs must already be installed.

For a native Qt/e-paper app that draws its own screen:
\begin{shell}
~/inkline shortcut register a --epaper /home/root/my-app
\end{shell}
One native app runs at a time, pausing Inkline and its shells. The launcher
supplies Qt e-paper settings and RAM caches, discards logs and resumes Inkline
when the app exits. Apps should use Qt input without exclusive keyboard grabs.

\begin{notice}
\key{Opt+RightAlt+T is permanent.} It cannot be registered or deregistered.
It stops a native shortcut app and returns to the existing terminals.
Do not hold Ctrl for this chord. Caps Lock remapping applies inside Inkline.
\end{notice}

\key{Emergency:} hold Opt+RightAlt+Backspace for two seconds to stop the shortcut
app and its descendants and restart Inkline. All terminals close; unsaved
work is lost. Root apps can defeat recovery by grabbing input, changing
services or exhausting resources. Restarting the device is the final fallback.

\topic{Install, update, keep the manual}
The bundle targets \key{reMarkable 2, firmware 3.27}, using the stock Qt 6.8
e-paper libraries. Current downloads and checksum instructions are at:

\InstallLink

The installer checks hardware, firmware, fonts, Qt startup, RAM temporary
directories and disabled swap. Releases live under
\path{/home/root/.local/share/inkline}; the keyboard launcher starts at boot.
Updates preserve open terminals. \key{Quit and reopen Inkline to use the
new version.} Previous releases remain for rollback and occupy storage.

This single PDF contains all nine languages and ships in the bundle.
Installation attempts to import it into \key{My files} through the USB web
interface when that interface is enabled and notebooks are running.
A successful import is recorded to avoid importing the same manual twice.

If automatic import is unavailable, connect USB, quit Inkline and enable
\key{USB web interface} in the notebook storage settings. From SSH, run:
\begin{shell}
~/inkline manual
\end{shell}
Alternatively, download this PDF and import it through the USB browser
interface or reMarkable's desktop/mobile app. The importer never edits the
notebook database directly or interrupts terminals. An imported manual stays
in My files after Inkline is uninstalled.

From SSH, return to notebooks or remove Inkline with:
\begin{shell}
~/inkline stop
~/inkline uninstall
\end{shell}
Optional software, including the recommended LaTeX installation, is documented
on the website.

\topic{How it works; storage and RAM}
Qt delivers keyboard input to a shared mapper for shortcuts and Caps Lock.
libghostty-vt encodes terminal keys. Each slot connects its program to a
separate pseudo-terminal (PTY). Shells receive \code{HOME=/home/root};
interactive Bash reads \code{~/.bashrc}.

Output passes through the terminal parser and sixel decoder. libghostty
maintains cells, cursor state, alternate screens, history and kitty images.
The renderer keeps a grayscale image and redraws changed rows; Qt's stock
e-paper backend updates the panel. An independent evdev daemon handles
global shortcuts and the power key; systemd handles suspend and wake.

Each terminal caps core allocations at 64 MiB, kitty images at 16 MiB per
screen and retained sixel images at 16 MiB. Rendering and app memory are
additional. Empty slots allocate no terminal. Nine large apps can exceed RAM.

Temporary graphics, caches and session state use RAM. Installed software,
saved documents, settings saved at exit and app-written files use persistent
storage. Inkline does not prevent arbitrary programs from writing files.

The optional Python package is stock Python. To reduce cache writes, put
these preferences in the tablet's \code{~/.bashrc}:
\begin{shell}
export PYTHONDONTWRITEBYTECODE=1
export PIP_NO_CACHE_DIR=1
\end{shell}
Run \code{source ~/.bashrc} or open a new shell. Use
\code{python3 -m pip install --no-compile PACKAGE} to avoid pip's separate
install-time bytecode compilation. Python isolated mode ignores Python
environment variables. See the website for the full Python guide.

Inkline is GPL-3.0-or-later software. Bundles include matching source archives
and dependency notices. Test records distinguish automated checks from
physical behaviors that still need confirmation.
